\section{An analytic proof of the transition}\label{app:analytic-threshold}

To prove Theorem~\ref{thm:threshold}, we turn the sign of the central
curvature into a comparison of positive weights. The moment equation
of Gu\'erin, Laulin, Raschel and Simon~\cite[(15)]{glrs2025}
supplies the integral representation.

Put
\begin{equation}\label{eq:an-H}
 \chi_k=4^{-k}\binom{2k}{k},\qquad
 H_a(y)=1-\sum_{k\ge1}\frac{\chi_k}{2ak-1}y^k.
\end{equation}
The function $H_a$ decreases strictly on $[0,1)$.
For $a=1$ it is $\sqrt{1-y}$. For $a<1$ every subtracted
coefficient is larger, so $H_a(1)<0$. The series at one converges,
since its coefficients are bounded by a constant depending on $a$
times $\chi_k/(2k-1)$, whose sum is one.
There is consequently a unique root $Y_a\in(0,1)$.
The coefficients also give
\begin{equation}\label{eq:an-H-ode}
 H_a(y)-2ayH_a'(y)=(1-y)^{-1/2}.
\end{equation}
Since $\partial_aH_a(y)>0$ for $y>0$ and $H_a'(Y_a)<0$,
the root is smooth and $Y_a'>0$.

\begin{lemma}\label{lem:curvature-weight}
For $1/2<a<1$,
\begin{equation}\label{eq:an-curvature}
 g_a''(0)=-\frac{1}{\pi\Gamma(3-3a)}
 \int_0^{Y_a}y^{1/2}H_a(y)^{2-3a}
             \bigl[1-(1+a)^2y\bigr]\,d y.
\end{equation}
The integral is absolutely convergent.
\end{lemma}

\begin{proof}
We multiply the limit by an independent exponential power to cancel
the gamma factors in its moments.
Let $E$ be an independent exponential random variable of mean one
and set $Z_a=E^a\mathcal L_a$. With the normalization of~\cite{glrs2025},
\[
 m_n=\frac{\Gamma(1+an)}{n!}\E \mathcal L_a^n,
 \qquad m_0=m_1=1.
\]
The bound $0<m_n\le C_a^n$ follows from
\cite[Lemma~2.10 and Corollary~2.14]{glr2026}. Applied to
the even moments, it gives $\E\exp(\delta|Z_a|)<\infty$ for small
$\delta>0$. Hence $M_a(z)=\E e^{zZ_a}$ is analytic on a strip
about the imaginary axis. Its power series and the moment equation give
\begin{equation}\label{eq:an-moment-ode}
 M_a(z)+azM_a'(z)
 =\frac{1+a}{2}M_a(z)^2
  +\frac{1-a}{2}M_a(z)M_a(-z).
\end{equation}
The identity holds on the strip by analytic continuation.
If $M_a(it)=A(t)+iB(t)$, it becomes
\begin{equation}\label{eq:an-AB}
 atA'=A^2-A-aB^2,\qquad
 atB'=((1+a)A-1)B.
\end{equation}

For $0<y<Y_a$ define
\begin{equation}\label{eq:an-param}
 t=\sqrt y\,H_a(y)^{-a},\qquad
 A(t)=H_a(y)\sqrt{1-y},\qquad B(t)=H_a(y)\sqrt y.
\end{equation}
This is the imaginary-axis form of the closed expression in
\cite[Proposition~3.2]{glrs2025}, with their $G(it)=i\sqrt y$
and $(G(it)/(it))^{1/a}=H_a(y)$. We verify the continuation
directly from~\eqref{eq:an-AB}.
The map from $y$ to $t$ increases from zero to infinity.
Equation~\eqref{eq:an-H-ode} gives
\[
 \frac{1}{t}\frac{dt}{dy}
 =\frac{1}{2yH_a(y)\sqrt{1-y}},
\]
and substitution shows that~\eqref{eq:an-param} satisfies~\eqref{eq:an-AB}.
Since $H_a(0)=1$, the analytic map $y\mapsto yH_a(y)^{-2a}$
has derivative one at zero. The identity $t^2=yH_a(y)^{-2a}$
therefore makes $y$ analytic in $t^2$ near zero.
Thus $A$ is even and analytic, while $B=tH_a(y)^{1+a}$ is odd
and analytic, with $A(0)=1$ and $B'(0)=1$.
These data determine both power series. The coefficient
recurrence has denominator $na-1$ for even $n\ge2$, and $na-a$
for odd $n\ge3$, both positive. Thus the parametrization agrees locally with
$M_a(it)$. Ordinary uniqueness for~\eqref{eq:an-AB} on $t>0$
extends the agreement to every positive $t$.

Write $f(t)=\operatorname{Re}\E e^{it\mathcal L_a}$. Exponential randomization gives
\[
 A(t)=\int_0^\infty e^{-v}f(tv^a)\,d v.
\]
We use the Fourier decay in Lemma~\ref{lem:analytic-density},
which in particular makes every polynomial multiple of $f$
integrable at infinity. For $0<\Re s<1/a$, absolute Fubini gives
\begin{equation}\label{eq:an-mellin}
 J(s):=\int_0^\infty t^{s-1}A(t)\,d t
 =\Gamma(1-as)\mathcal F(s),\qquad
 \mathcal F(s)=\int_0^\infty t^{s-1}f(t)\,d t.
\end{equation}
The function $\mathcal F$ is holomorphic on $\Re s>0$.
The integral defining $J$ diverges at $s=3$, the value needed for curvature.
Two integrations by parts give a convergent expression there.

We send the root to zero before integrating by parts. Set $\theta=\arcsin\sqrt y$ and
\[
 u=F(\theta)=\sin^{-1/a}\theta\,H_a(\sin^2\theta),\qquad
 w(u)=a\cos\theta\sin^{1/a}\theta.
\]
Here $\theta$ decreases from $\arcsin\sqrt{Y_a}$ to zero as
$u$ increases from zero to infinity. From~\eqref{eq:an-H-ode},
\[
 F'(\theta)=-\frac1a\sin^{-1-1/a}\theta,\qquad t=u^{-a}.
\]
Changing variables in $J$ gives
$J(s)=\int_0^\infty u^{-as}w(u)\,d u$.
Since $F'(\arcsin\sqrt{Y_a})\ne0$, the function $w$ is smooth
at zero. At infinity, $\theta\sim u^{-a}$, so its definition
and the derivative formulas below give
$w(u)=O(u^{-1})$, $w'(u)=O(u^{-2})$ and $w''(u)=O(u^{-3})$.
For $0<\Re s<1/a$, the boundary terms therefore vanish twice,
and~\eqref{eq:an-mellin} becomes
\begin{equation}\label{eq:an-mellin-continued}
 \mathcal F(s)=\frac{1}{\Gamma(3-as)}
       \int_0^\infty u^{2-as}w''(u)\,d u.
\end{equation}
The right side is holomorphic on $0<\Re s<3/a$ by the bounds
at infinity and boundedness of $w''$ at zero.
Thus~\eqref{eq:an-mellin-continued} holds throughout this larger
strip by the identity theorem, and $s=3$ is permitted.

Differentiating through $u=F(\theta)$ gives
\[
 \begin{split}
 w'(u)&=-a\sin^{2/a}\theta
                  [1-(1+a)\sin^2\theta],\\
 w''(u)&=2a\sin^{3/a}\theta\cos\theta
                  [1-(1+a)^2\sin^2\theta].
 \end{split}
\]
Insert the second expression in~\eqref{eq:an-mellin-continued}
at $s=3$ and return to $y=\sin^2\theta$.
The powers of $y$ reduce to $y^{1/2}$, giving the integral in
\eqref{eq:an-curvature}. Fourier inversion supplies the factor
$g_a''(0)=-\mathcal F(3)/\pi$.
Finally, $H_a$ has a simple zero at $Y_a$, and $2-3a>-1$,
so the endpoint singularity is integrable.
\end{proof}

\begin{proof}[Proof of Theorem~\ref{thm:threshold}]
\emph{Signs outside a short interval.}

Let $\mu_a$ be the probability law on $(0,Y_a)$ with density proportional to
$y^{1/2}H_a(y)^{2-3a}$, and put
\begin{equation}\label{eq:an-score}
 \mathcal R(a)=(1+a)^2\int y\,d\mu_a(y).
\end{equation}
Lemma~\ref{lem:curvature-weight} says that $g_a''(0)$ has the sign
of $\mathcal R(a)-1$.

When $a\le2/3$, the extra weight $H_a^{2-3a}$ decreases in $y$.
By Chebyshev's association inequality, it lowers the mean of the density proportional to
$y^{1/2}$, giving $\int y\,d\mu_a(y)\le3Y_a/5$.
At $a=2/3$, the first three terms of~\eqref{eq:an-H} give
\[
 H_{2/3}(3/5)
 \le1-\frac9{10}-\frac{81}{1000}-\frac9{400}
 =-\frac7{2000}<0.
\]
Since $Y_a$ increases, $Y_a\le Y_{2/3}<3/5$ for $a\le2/3$.
Thus $\mathcal R(a)<(5/3)^2(3/5)^2=1$.

When $a\ge2/3$, the extra weight increases, giving a mean of at least
$3Y_a/5$. We need the bound
$Y_{2/3}>117/200$. To prove it, note that $\chi_k$ decreases.
At $y=117/200$, the terms with $k\ge4$ are at most
\[
 \sum_{k\ge4}\frac{\chi_ky^k}{4k/3-1}
 \le\frac{\chi_4y^4}{(13/3)(1-y)}.
\]
Together with the first three terms this gives the exact bound
\[
 H_{2/3}(117/200)\ge\frac{929771}{135987200}>0.
\]
For every $a\ge69/100$, therefore,
\[
 \mathcal R(a)>
 \left(\frac{169}{100}\right)^2\frac35\frac{117}{200}
 =\frac{10024911}{10000000}>1.
\]
Any zero must lie between $2/3$ and $69/100$.

\emph{A single crossing.}

We now work on $I=[2/3,69/100]$. To compare densities on a fixed
interval, rescale by $y=Y_au$, so $u$ always lies in $(0,1)$.
We separate the endpoint zero by putting
\begin{equation}\label{eq:an-weights}
 w_k(a)=\frac{\chi_kY_a^k}{2ak-1},\qquad
 K_a(u)=\frac{H_a(Y_au)}{1-u}
       =\sum_{k\ge1}w_k(a)(1+u+\cdots+u^{k-1}).
\end{equation}
Since $\sum w_k=1$, we have $K_a\ge1$ and $\partial_uK_a\ge0$.
Writing $s_a=3a-2\in[0,7/100]$, the rescaled measure has density
proportional to
\begin{equation}\label{eq:an-rescaled-weight}
 W_a(u)=u^{1/2}(1-u)^{-s_a}K_a(u)^{-s_a}.
\end{equation}
To make the mean increase, we show that the score
$\partial_a\log W_a(u)$ is strictly increasing in $u$ and use
the covariance identity for the derivative of an expectation.

First, $Y_a<2/3$ on $I$. The first three terms give
\[
 H_{69/100}(2/3)\le-\frac{9767}{7087608}<0,
\]
and $Y_a$ increases. All series and derivatives used below
therefore converge uniformly on $I\times[0,1]$.
Define the falling-factorial moments
\[
 M_1=\sum k w_k,\qquad M_2=\sum k(k-1)w_k,\qquad
 M_3=\sum k(k-1)(k-2)w_k.
\]
For $G(y)=(1-y)^{-1/2}$, equation~\eqref{eq:an-H-ode} and its
first two derivatives at $Y=Y_a$ give
\begin{equation}\label{eq:an-factorial-moments}
 \begin{split}
 2aM_1&=G(Y),\\
 2aM_2&=YG'(Y)-(2a-1)M_1,\\
 2a(M_3+2M_2)&=Y^2G''(Y)+M_2.
 \end{split}
\end{equation}
In particular,
\begin{equation}\label{eq:an-Ku}
 0\le\partial_uK_a(u)\le\frac{M_2}{2}
 \le\frac{YG'(Y)}{4a}
 \le\frac{3\sqrt3}{8}<\frac{21}{32}.
\end{equation}

We next control memory derivatives of the kernel. Write
$d_k=2k/(2ak-1)$. Differentiating $\sum w_k=1$ gives
\begin{equation}\label{eq:an-root-derivative}
 \frac{Y_a'}{Y_a}M_1=\sum d_kw_k\le6,
 \qquad \frac{Y_a'}{Y_a}\le\frac{12a}{G(Y)}.
\end{equation}
Here $d_k$ decreases in $k$ and $d_1\le6$ on $I$.
The logarithmic derivative $\partial_a\log w_k=kY_a'/Y_a-d_k$
strictly increases in $k$. Since $1+u+\cdots+u^{k-1}$
is nondecreasing in $k$, the covariance identity for
derivatives of expectations gives $\partial_aK_a(u)\ge0$.
The probability law here is formed by the $w_k$.

In $\partial_a\partial_uK_a$, discard the negative term
$-d_kw_k$ from $\partial_aw_k$, then bound the derivative
of $1+u+\cdots+u^{k-1}$ by $k(k-1)/2$. Equations
\eqref{eq:an-factorial-moments} and~\eqref{eq:an-root-derivative} give
\begin{align}
 \partial_a\partial_uK_a
 &\le\frac{Y_a'}{Y_a}\frac{M_3+2M_2}{2}\notag\\
 &\le3\left(\frac{M_2}{G(Y)}+
                       \frac{Y^2G''(Y)}{G(Y)}\right)
 \le3\left(\frac34+3\right)=\frac{45}{4}.
 \label{eq:an-Kau}
\end{align}
The last bounds use
$M_2/G(Y)\le Y/[4a(1-Y)]$ and
$Y^2G''(Y)/G(Y)=3Y^2/[4(1-Y)^2]$.
Since $K_a\ge1$, $\partial_aK_a\ge0$ and
$\partial_uK_a\ge0$, we obtain
\[
 \partial_a\left(\frac{\partial_uK_a}{K_a}\right)
 =\frac{\partial_a\partial_uK_a}{K_a}
  -\frac{\partial_uK_a\,\partial_aK_a}{K_a^2}
 \le\frac{45}{4}.
\]
In~\eqref{eq:an-rescaled-weight}, the endpoint term dominates the kernel
terms, giving the strict bound
\begin{align}
 \partial_u\partial_a\log W_a(u)
 &=\frac3{1-u}-3\frac{\partial_uK_a}{K_a}
    -s_a\partial_a\left(\frac{\partial_uK_a}{K_a}\right)\notag\\
 &\ge3-\frac{63}{32}-\frac7{100}\frac{45}{4}
   =\frac{39}{160}>0.
 \label{eq:an-increasing-score}
\end{align}
The score $\partial_a\log W_a$ is therefore strictly increasing
in $u$. If $\nu_a$ denotes the normalized measure in
\eqref{eq:an-rescaled-weight}, differentiation gives
\[
 \frac{d}{da}\int u\,d\nu_a(u)
 =\operatorname{Cov}_{\nu_a}\bigl(u,\partial_a\log W_a(u)\bigr)>0.
\]
This is the strict form of Chebyshev's association inequality.
Indeed, the covariance is
one half the expectation of
$(U-V)(S(U)-S(V))$ for independent $U,V\sim\nu_a$ and the
strictly increasing score $S$. Differentiation under the
integral is justified by the integrable bound
$C(1-u)^{-7/100}(1+|\log(1-u)|)$ on this compact parameter interval.

It follows that
\[
 \mathcal R(a)=(1+a)^2Y_a\int u\,d\nu_a(u)
\]
has strictly positive derivative in the interior of $I$. The signs already proved
and the intermediate value theorem give
exactly one zero of $\mathcal R-1$. There are no zeros outside $I$.
Finally,~\eqref{eq:an-curvature} writes
\[
 g_a''(0)=\frac{\mathcal Z_a}{\pi\Gamma(3-3a)}\,[\mathcal R(a)-1],
 \qquad \mathcal Z_a=\int_0^{Y_a}y^{1/2}H_a(y)^{2-3a}\,d y>0.
\]
The normalizing factor is continuously differentiable near the
zero, by the same rescaling and domination. At the zero its
derivative contributes nothing, so
$\partial_ag_a''(0)=\mathcal Z_a\mathcal R'(a)/[\pi\Gamma(3-3a)]>0$.
The shape alternatives in Theorem~\ref{thm:two-modes} now give
unimodality below $a_*$ and bimodality above it. At the critical
point, let $a\uparrow a_*$ from below. Joint continuity preserves
$g_a(x)\ge g_a(y)$ for $0\le x<y$. Thus the critical density is
nonincreasing on the positive half-line too. Analyticity excludes
a constant interval, so its unique mode is at zero.
\end{proof}

The root bounds also determine the sign of the fourth derivative
at the threshold.

\begin{lemma}\label{lem:analytic-quartic}
For $2/3\le a\le69/100$, we have $g_a^{(4)}(0)<0$.
In particular, $B_*=-g_{a_*}^{(4)}(0)>0$.
\end{lemma}

\begin{proof}
We first extend the Mellin calculation by one integration by parts.
With $u,w$ as in the proof of Lemma~\ref{lem:curvature-weight}, put
\[
 c_a=(1+a)^2,\qquad
 P_a(y)=\frac3a(1-y)(1-c_ay)-y(1-c_ay)-2c_ay(1-y).
\]
Since $dy/du=-2ay^{1+1/(2a)}\sqrt{1-y}$, differentiation gives
$w'''(u)=-2a^2y^{2/a}P_a(y)$.
This derivative is bounded at zero and is $O(u^{-4})$ at infinity.
The boundary terms vanish for $0<\Re s<3/a$, so
\eqref{eq:an-mellin-continued} gives
\[
 \mathcal F(s)=-\frac1{\Gamma(4-as)}
                \int_0^\infty u^{3-as}w'''(u)\,d u.
\]
The right side is holomorphic for $0<\Re s<4/a$.
For $1/2<a<4/5$, we may therefore set $s=5$ and use
$g_a^{(4)}(0)=\mathcal F(5)/\pi$ to obtain
\begin{equation}\label{eq:an-fourth-derivative}
 g_a^{(4)}(0)=\frac{a}{\pi\Gamma(4-5a)}
 \int_0^{Y_a}y^{3/2}(1-y)^{-1/2}H_a(y)^{3-5a}P_a(y)\,d y.
\end{equation}
The integral converges absolutely because $H_a$ has a simple zero
and $3-5a>-1$.

Now let $a\in[2/3,69/100]$. The root bounds in the preceding proof
give $1/2<Y_a<2/3$. For the coefficient comparison only, we use
the larger interval $2/3\le a\le7/10$. Expanding $P_a$ gives, for $y\ge0$,
\[
 \begin{split}
 P_a(y)&=\frac3a-\left(2a^2+7a+9+\frac6a\right)y
                    +\left(3a^2+9a+9+\frac3a\right)y^2\\
       &\le\frac92-23y+22y^2=:Q(y).
 \end{split}
\]
Here termwise lower bounds give
$8/9+14/3+9+60/7=1457/63>23$ for the coefficient of $-y$,
and termwise upper bounds give
$147/100+63/10+9+9/2=2127/100<22$ for that of $y^2$.
Since $Q(0)>0$, $Q(1/2)=-3/2$ and $Q(2/3)=-19/18$,
this quadratic has exactly one root $r$ in $(0,Y_a)$.
It is positive before $r$ and negative after it.
For $1/2\le Y\le2/3$, direct integration gives
\[
 Y^{-5/2}\int_0^Y y^{3/2}Q(y)\,d y
 =\frac95-\frac{46}{7}Y+\frac{44}{9}Y^2
 \le-\frac{83}{315}<0,
\]
where convexity reduces the bound to the two endpoints.

The remaining weight $R_a(y)=(1-y)^{-1/2}H_a(y)^{3-5a}$
is positive and increasing on $(0,Y_a)$.
The sign change at $r$ therefore gives
$R_a(y)Q(y)\le R_a(r)Q(y)$ throughout this interval. Consequently,
\[
 \int_0^{Y_a}y^{3/2}R_a(y)P_a(y)\,d y
 \le R_a(r)\int_0^{Y_a}y^{3/2}Q(y)\,d y<0.
\]
The prefactor in~\eqref{eq:an-fourth-derivative} is positive on
$[2/3,69/100]$, which proves the claim.
\end{proof}

